\chapter{Fraunhofer Diffraction}

\section*{Introduction}

For a single slit of width $a$, the Fraunhofer diffraction pattern has intensity $I(\theta) = I_0\,\text{sinc}^2(\beta)$, where $\beta = \pi a \sin\theta/\lambda$ and $\text{sinc}(x) = \sin(x)/x$. The central maximum is the brightest, and the pattern consists of a series of diminishing side lobes separated by zeros at $\sin\theta = n\lambda/a$ (for $n = \pm 1, \pm 2, \ldots$). For a double slit, the pattern is modulated by an interference envelope, producing the classic Young's fringes. For a diffraction grating with $N$ slits, the principal maxima are $N$ times narrower than for a single slit, making gratings ideal for spectroscopy. The Fraunhofer approximation is valid when the observation distance $D \gg a^2/\lambda$---physically, when the wavefronts arriving at the aperture are effectively planar.


\[
I(\theta) = I_0\,\text{sinc}^2(\beta),\qquad \beta = \frac{\pi a \sin\theta}{\lambda}
\]

\section*{Fundamental Equations Used}

The derivation starts from the Huygens--Fresnel principle and the Rayleigh--Sommerfeld diffraction integral in the far-field limit.

\section*{Assumptions}

\textbf{Assumptions:} The observation point is in the far field (Fraunhofer regime): $D \gg a^2/\lambda$. The aperture is illuminated by a coherent, monochromatic plane wave. The aperture is in a thin, opaque screen. Diffraction angles are small ($\sin\theta \approx \theta$).


\textbf{Limitations:} Fails in the near field (Fresnel diffraction), where the curvature of the wavefront across the aperture matters. Does not account for vector effects (polarization-dependent diffraction) unless a vector diffraction theory is used. For apertures comparable to $\lambda$, scalar diffraction theory breaks down and full electromagnetic solutions are needed.


\section*{Mathematical Derivation}

\textbf{Step 1: Huygens--Fresnel principle and the diffraction integral.}

According to the Huygens--Fresnel principle, every point on a wavefront acts as a secondary source of spherical wavelets. The field at an observation point $P$ due to an aperture in the $xy$-plane is given by the Rayleigh--Sommerfeld diffraction integral. In the Fraunhofer (far-field) limit, where the observation distance $D$ satisfies $D \gg a^2/\lambda$ ($a$ is the aperture size), the field simplifies to a Fourier transform:
\begin{align}
U(x', y') = \frac{e^{ikD}e^{ik(x'^2+y'^2)/(2D)}}{i\lambda D}\iint_{\text{aperture}} U_0(x,y)\,e^{-ik(xx'+yy')/D}\,dx\,dy.
\end{align}

\textbf{Step 2: Single-slit geometry.}

Consider a slit of width $a$ along the $x$-axis, infinite in the $y$-direction, illuminated by a plane wave of amplitude $U_0$. The aperture function is $U_0(x) = 1$ for $|x| \leq a/2$ and zero otherwise. The far-field pattern depends only on $x'$. Using the small-angle approximation ($\sin\theta \approx x'/D$) and defining $k_x = k\sin\theta = 2\pi\sin\theta/\lambda$, the diffracted field becomes
\begin{align}
U(\theta) \propto \int_{-a/2}^{a/2} e^{-ik_x x}\,dx.
\end{align}

\textbf{Step 3: Evaluate the integral.}

Performing the integration:
\begin{align}
U(\theta) &\propto \left[\frac{e^{-ik_x x}}{-ik_x}\right]_{-a/2}^{a/2} = \frac{e^{-ik_x a/2} - e^{ik_x a/2}}{-ik_x} = \frac{2\sin(k_x a/2)}{k_x}.
\end{align}
Substituting $k_x = 2\pi\sin\theta/\lambda$:
\begin{align}
U(\theta) &\propto a\,\frac{\sin(\pi a\sin\theta/\lambda)}{\pi a\sin\theta/\lambda} = a\,\text{sinc}(\beta),
\end{align}
where $\beta = \pi a\sin\theta/\lambda$ and $\text{sinc}(x) = \sin(x)/x$.

\textbf{Step 4: Intensity pattern.}

The intensity is $I(\theta) = |U(\theta)|^2$. Therefore:
\begin{align}
I(\theta) = I_0\,\text{sinc}^2(\beta) = I_0\left(\frac{\sin\beta}{\beta}\right)^{\!2}, \qquad \beta = \frac{\pi a\sin\theta}{\lambda},
\end{align}
where $I_0$ is the peak intensity at $\theta = 0$ (the central maximum), obtained by setting $\beta = 0$ so that $\text{sinc}(0) = 1$.

\textbf{Step 5: Minima and physical interpretation.}

Dark fringes occur when $\sin\beta = 0$ but $\beta \neq 0$, i.e., $\beta = n\pi$ for integer $n \neq 0$. This gives $\sin\theta_n = n\lambda/a$, confirming the zeros of the pattern. The central maximum spans from the first minimum on one side ($n = -1$) to the other ($n = +1$), giving an angular width $\Delta\theta \approx 2\lambda/a$---twice the width of each side lobe. This inverse relationship between slit width and pattern width is a direct manifestation of the Fourier uncertainty principle: confining the wave in real space (narrow slit) spreads it in frequency space (wide diffraction pattern).

\section*{Characteristics of the Equation}

The central maximum of the Fraunhofer pattern has width $\Delta\theta \approx 2\lambda/a$ (from first zero to first zero), so narrower apertures produce wider patterns---a manifestation of the uncertainty principle. The side lobes decrease in intensity as $1/\beta^2$ for large $\beta$. For a circular aperture, the pattern becomes the Airy disk: $I(\theta) = I_0\left[2J_1(x)/x\right]^2$ with $x = \pi D \sin\theta/\lambda$, and the first zero occurs at $\sin\theta = 1.22\lambda/D$. The total power transmitted is conserved: narrowing the slit concentrates the same power into a wider pattern.
\section*{Solution Methods}

\textbf{Fourier transform method:} Compute the Fourier transform of the aperture function and square its magnitude. This is the most general approach. \textbf{Phasor addition:} Divide the aperture into Huygens wavelets and sum their contributions with appropriate phase differences. For a single slit, this yields the sinc function. \textbf{Numerical computation:} For complex apertures (arbitrary shapes, apodized gratings), compute the Fourier transform numerically using FFT algorithms. \textbf{Rayleigh--Sommerfeld integral:} The exact diffraction integral can be evaluated numerically when the Fraunhofer approximation is not valid.
\section*{Verifications}

For a single slit, verify that the first zero occurs at $\sin\theta = \lambda/a$ by setting $\beta = \pi$. Check that the central maximum is twice as wide as the side lobes. For a double slit with separation $d$, verify that the interference pattern is modulated by the single-slit envelope: $I \propto \text{sinc}^2(\beta)\cos^2(\delta/2)$ with $\delta = 2\pi d\sin\theta/\lambda$. For a grating with $N$ slits, verify that the principal maxima are $N$ times narrower and $N^2$ times brighter than the single-slit pattern. Compare numerical FFT calculations with analytical results for rectangular and circular apertures.
\section*{Applications}

Fraunhofer diffraction determines the resolving power of telescopes, microscopes, and cameras---the Rayleigh criterion ($\theta_{\min} = 1.22\lambda/D$) comes directly from the Airy disk. Diffraction gratings use the Fraunhofer pattern to separate spectral lines, enabling spectroscopy in astronomy, chemistry, and materials science. X-ray diffraction (Bragg's law) uses Fraunhofer diffraction from crystal lattices to determine atomic structure. In telecommunications, diffraction limits the focusing of laser beams through fibers. In photolithography, diffraction limits the minimum feature size that can be patterned on semiconductor wafers.
\section*{What Richard Feynman Would Have Asked}

\subsection*{1. Plain-English Translation}
\textit{``What is diffraction really? Not the math---what is the physics?''}

Diffraction is what happens when a wave tries to squeeze through a hole that is not much bigger than its wavelength. The wave cannot simply pass through as a narrow beam---it spreads out. The narrower the slit, the more it spreads. This is not a failure of the wave; it is an inevitable consequence of the wave nature of light. A wave confined to a narrow opening must develop a spread in the transverse direction, just as a particle confined to a small region must develop a spread in momentum (Heisenberg).

The sinc-squared pattern is nature's fingerprint of the rectangular aperture. Every shape produces its own characteristic diffraction pattern---a circle produces the Airy disk (Bessel function), a square produces a sinc-squared in both directions, and a random shape produces a random-looking pattern. The pattern is always the Fourier transform of the aperture shape, because diffraction is fundamentally about decomposing the wave into plane waves traveling in different directions.

\subsection*{2. Localized Mechanics}
\textit{``At a single point on the screen, how do the waves from different parts of the slit combine?''}

At the center of the pattern, all wavelets from across the slit arrive in phase---they all travel the same distance. They add constructively to produce maximum intensity. As you move off-center, wavelets from opposite edges of the slit travel different distances. At the first dark fringe, the path difference between the top and bottom of the slit is exactly one wavelength. The wavelets from the upper half of the slit are exactly out of phase with those from the lower half, and they cancel pairwise.

The side lobes arise because the cancellation is imperfect. In the first side lobe, two-thirds of the slit cancels and one-third contributes---hence the much lower intensity (about $4.7\%$ of the peak). The second side lobe has even less uncancelled area. Each successive side lobe corresponds to a more finely balanced cancellation across the slit.

\subsection*{3. Testing Extreme Limits}
\textit{``What happens when the slit is much larger or much smaller than the wavelength?''}

Wide Slit ($a \gg \lambda$): The diffraction pattern is confined to very small angles---the light passes through essentially as a straight beam. This is the geometric optics limit. The central maximum is so narrow that you see a sharp shadow of the slit.

Narrow Slit ($a \sim \lambda$): The pattern spreads over nearly the entire forward hemisphere. The concept of a ``beam'' breaks down---light scatters in all forward directions. The sinc function is still valid, but its first zero is at $\sin\theta = 1$, meaning the pattern extends to $90^\circ$.

Single Atom ($a \sim 0$): The pattern becomes isotropic---the atom radiates equally in all directions (Rayleigh scattering). The atom is a point source, and its diffraction pattern is uniform. This is why atoms scatter light in all directions, making the sky blue.

\subsection*{4. Dimensionality and Geometry}
\textit{``What changes if we go from a slit to a circular hole?''}

The mathematical structure changes because the symmetry changes. A slit has rectangular symmetry, so the pattern involves $\text{sinc}$ functions (one-dimensional Fourier transforms). A circular aperture has circular symmetry, so the pattern involves Bessel functions (the Fourier transform of a circular function). The Airy disk---the circular diffraction pattern---has the same qualitative features as the single-slit pattern (central maximum, diminishing rings) but different quantitative details.

In one dimension (a slit), the side lobes decay as $1/\beta^2$. In two dimensions (a circular aperture), they decay faster because the area of each ring increases. The circular symmetry ``averages out'' the oscillations more efficiently. This is why circular apertures produce cleaner images than rectangular ones---the side lobes contribute less to image degradation.

\subsection*{5. Cross-Disciplinary Connection}
\textit{``Where else does this diffraction pattern appear outside of optics?''}

Diffraction is a universal wave phenomenon. Sound diffracting through a doorway produces the same sinc-squared pattern. Radio waves diffracting around buildings produce the same interference patterns. Neutron diffraction reveals crystal structures the same way X-rays do. Even matter waves (electrons, neutrons, fullerenes) diffract when passed through gratings---de Broglie's hypothesis confirmed by experiment.

In signal processing, the Fourier transform relationship is the foundation of frequency analysis. The diffraction pattern of an aperture is mathematically identical to the frequency spectrum of a rectangular pulse---the sinc function appears in both. In medical imaging (CT, MRI), the reconstruction algorithms rely on the same Fourier relationship between the measured data and the image.

%=============================================================
\section*{FAQ}

\textbf{Q: Why is the pattern called ``far-field'' diffraction?}

A: ``Far-field'' means the observation distance is much greater than $a^2/\lambda$. At this distance, the wavefronts arriving at the aperture from any point source are effectively planar, and the diffraction pattern is the Fourier transform of the aperture. In the near field (Fresnel region), the wavefront curvature matters and the pattern is more complex.

\textbf{Q: How does a diffraction grating differ from a single slit?}

A: A grating has many equally spaced slits. The interference between slits produces sharp principal maxima at angles $\sin\theta = n\lambda/d$ (where $d$ is the slit spacing), while the single-slit diffraction provides an envelope. The more slits, the sharper the peaks---this is the basis of high-resolution spectroscopy.

\textbf{Q: Can you beat the diffraction limit?}

A: Not with far-field optics using propagating waves. However, near-field techniques (SNOM/NSOM) use evanescent waves to achieve sub-wavelength resolution. Super-resolution microscopy techniques (STED, PALM, STORM) use fluorescent markers and clever illumination to bypass the classical diffraction limit.
\section*{Important Bibliography}

\begin{enumerate}
\item Hecht, E., \textit{Optics}, 5th ed. (McGraw-Hill, 2017), Chapter 10.
\item Born, M. and Wolf, E., \textit{Principles of Optics}, 7th ed. (Cambridge University Press, 1999), Chapter 8.
\item Goodman, J.W., \textit{Introduction to Fourier Optics}, 4th ed. (W.H.\ Freeman, 2017), Chapters 4--5.
\item Fraunhofer, J., ``Bestimmung des Brechungs- und Farbenverm\"ogens des Glases,'' \textit{Annalen der Physik} \textbf{74}, 271--300 (1817).
\end{enumerate}

